\documentclass{article}


% if you need to pass options to natbib, use, e.g.:
%     \PassOptionsToPackage{numbers, compress}{natbib}
% before loading neurips_2022


% ready for submission
\usepackage{neurips_2022}


% to compile a preprint version, e.g., for submission to arXiv, add add the
% [preprint] option:
%     \usepackage[preprint]{neurips_2022}


% to compile a camera-ready version, add the [final] option, e.g.:
%     \usepackage[final]{neurips_2022}


% to avoid loading the natbib package, add option nonatbib:
%    \usepackage[nonatbib]{neurips_2022}


\usepackage[utf8]{inputenc} % allow utf-8 input
\usepackage[T1]{fontenc}    % use 8-bit T1 fonts
\usepackage{hyperref}       % hyperlinks
\usepackage{url}            % simple URL typesetting
\usepackage{booktabs}       % professional-quality tables
\usepackage{amsfonts}       % blackboard math symbols
\usepackage{nicefrac}       % compact symbols for 1/2, etc.
\usepackage{microtype}      % microtypography
\usepackage{xcolor}         % colors

% Our packages
\usepackage{caption}
\usepackage{subcaption}
\usepackage{tikz}
\usepackage{pgfplots}
\DeclareUnicodeCharacter{2212}{−}
\usepgfplotslibrary{groupplots,dateplot}
%\usetikzlibrary{external}
%\tikzexternalize[prefix=tikz/]
\usetikzlibrary{patterns,shapes.arrows}
\pgfplotsset{compat=newest}
\usetikzlibrary{shapes,decorations,arrows,calc,arrows.meta,fit,positioning}
\usetikzlibrary{positioning,calc}
\usetikzlibrary{backgrounds,fit}
\usetikzlibrary{arrows.meta}
\usetikzlibrary{arrows}

\usepackage{svg}
\usepackage{listings}
\usepackage{amsmath,amssymb}
\DeclareMathOperator*{\E}{\mathbb{E}}
\usepackage[ruled,vlined]{algorithm2e}


\usepackage{todonotes}
\newcommand{\frank}[1]{\todo[inline,color=yellow]{Frank: #1}}
\newcommand{\sam}[1]{\todo[inline,color=green]{Sam: #1}}
\newcommand{\noah}[1]{\todo[inline,color=orange]{Noah: #1}}
\newcommand{\katha}[1]{\todo[inline,color=white]{Katha: #1}}


%\title{Probabilistic AutoML in a Single Forward Pass}
% Probabilistic
% Causal
% 
\title{Meta-Learning Fast Tabular AutoML for Small Data}

\begin{document}


We extend BNNs to represent a subclass of Structural Causal Models (SCMs~\citep{causalinferencelearningbook}).
We focus in specific on functional SCMs, that have variables of the form $X = \{x_1,\dots,x_n\}$, where each variable $x_i$ is set by a deterministic function $f_i$ and some independent noise $\eps_i$ as
\begin{align}
    %x_i = f_i(X_{\text{PA}_{\mathcal{G}}(i)}) + \eps_i,
\end{align}
where $\mathcal{G}$ is an underlying causal graph in which $\text{PA}_{\mathcal{G}}(i)}$ refers to the parents of the $i$th node.
underlying causal DAG $(V,E)$ in which the value for each source variable $v \in V$, all nodes without parents, is defined purely as noise $x = \$ non variable $v \in V$ is defined by a the form
That is we use SCMs as a language to describe our hypotheses, we are inspired by science showing causal mechanisms to be a strong prior in human reasoning as well \citep{johnson2014explanatory, wojtowicz2020probability}.
\end{document}